\section{模型的评价与推广}

\subsection{模型的优点与不足}

\subsubsection{主要优点}

\begin{enumerate}
  \item \textbf{物理机理与凸优化深度耦合，零数据先验且求解高效}：针对多目标高轨抵近巡游的混合整数非凸难题，本文摒弃了对特定数据集的手工先验与耗时数小时的盲目随机搜索，提出三维 Clohessy--Wiltshire 能量积分两脉冲物理状态动态规划（3D-CW-VND-SE-DP），将三维倾角/偏心率耦合与短时跨面代价直接嵌入毫秒级精确动态规划中，自动识别伴随飞越并跳出局部极小；连续层进一步将脉冲矢量、施加时刻与新星初始轨道状态统一纳入虚拟控制二阶锥序列凸优化（VC-SCP），在 24 核并行下 7.4~min 内即将问题二压至 159.33~m/s，将问题三压至 28.11~m/s。
  \item \textbf{观测约束的动态切平面凸化严格保可行且全链条可信}：可观测区域中距离上界与双圆锥约束天然为二阶锥，非凸距离下界采用动态切平面内近似，确保每个凸子问题的无松弛解在连续空间中严格可行；三问全部 61 次观测均在 0.05~s 加密网格（最大距离采样误差 $\le 2.03$~mm）与 ATK 4.2.0 独立数值积分下 100\% 通过复核。
  \item \textbf{参数灵敏度清晰，工程稳健性强}：通过对安全裕量 $m$、R2 初始偏心率上限 $e_{\max}$ 与半长轴偏离上限 $\delta a$ 的系统延拓分析，定量阐明了问题三中起主导作用的是半长轴偏离上限 $\delta a$ 而非偏心率上限，且主方案对各工程参数的灵敏度均较低。
\end{enumerate}

\subsubsection{存在的不足}

\begin{enumerate}
  \item \textbf{未计入高阶轨道摄动}：本文遵循题面设定采用中心引力场二体模型，未考虑地球非球形引力（$J_2$ 项及田谐项）、日月三体引力与太阳光压摄动。在 30 天任务周期内，上述摄动会在地球同步轨道上引起数十千米量级的相位漂移，实际工程执行时需在本文二体方案基础上加入轨道保持微调。
  \item \textbf{脉冲瞬时假设与有限推力差异}：本文将每次轨道机动理想化为瞬时速度增量，未考虑实际化学推力器或电推力器的有限燃烧时长与姿态机动建立时间；对于单次接近 24~m/s 的较大脉冲，有限推力重力损耗会带来少量额外燃料消耗。
\end{enumerate}

\subsection{模型的推广方向}

本文建立的“三维相对运动物理状态动态规划 + 全耦合虚拟控制序列凸优化”框架具有高度模块化结构，可直接向以下工程场景推广：

\begin{enumerate}
  \item \textbf{高精度摄动动力学热启动}：由于连续层 VC-SCP 的灵敏度矩阵 $\bm\Phi_{rv}$、$\bm\Phi_{rt}$ 与 $\bm\Phi_{rx}$ 均通过轨迹传播器有限差分求得，只需将二体开普勒传播器替换为包含 $J_2$、日月引力与太阳光压的高精度数值积分器，即可以本文二体最优解为热启动初值，在数轮序列凸优化内直接收敛至全摄动环境下的精确可执行轨道。
  \item \textbf{有限推力连续弧段转换}：对单次较大的脉冲机动，可在脉冲时刻附近展开为对称有限推力燃烧弧段，在式~\eqref{eq:socp} 中将瞬时脉冲替换为分段常值推力加速度矢量，同样保持二阶锥凸优化形式，无缝衔接高比冲电推进或常规液体推力器约束。
  \item \textbf{多星星座（$N \ge 3$）协同巡游网络}：将第 7 章的双星经度分组推广为多星空间分区与任务负载均衡模型，每颗新增观测星的初始六维轨道状态均可作为独立块变量嵌入统一序列凸优化框架，从而直接应用于大规模地球同步轨道空间态势感知星座的协同部署与常态化巡检任务规划。
\end{enumerate}
